\section{Conclusion and Limitations}
Our model shows strong performance under sparse-view constraints, specifically when handling between 3 and 12 views.
The model demonstrates the importance of accurate dense point cloud initialization.
We introduce a modified depth loss that enables correct scene generalization by reducing depth ambiguities without introducing artifacts in low confidence regions.
In addition, we introduce normal and depth warping loss terms that improve alignment with the ground-truth surface geometry. Finally, we relax certain assumptions from PGSR to allow robust optimization in sparse-view settings.

Our model faces limitations when dealing with large datasets, as processing many input views with $\pi^3$ consumes a large amount of GPU memory, which is infeasible on consumer hardware. Additional limitations come from inaccurate depth estimations in specific scenes, such as the leaves scene from the LLFF dataset~\cite{mildenhall2019llff}. Future improvements could include the joint optimization of the camera poses and the Gaussian scene, which would result in improved reconstruction quality. Furthermore, the integration of generative priors could enhance the model's ability to maintain photometric and geometric consistency across occluded or sparse areas.
\FloatBarrier